\chapter{World Regions and Cultures}
\label{ch:world-cultures}
\index{world cultures}\index{lore}

\begin{tcolorbox}[colback=blue!5!white, colframe=blue!75!black, title=Using This Compendium]
\textbf{You only need to read the region your character is from or where the campaign starts.} The GM will tell you when other regions become relevant.

\medskip
\noindent Each entry includes:
\begin{itemize}
  \item \textbf{Fantasy subgenre} – the kind of stories you can expect.
  \item \textbf{Attitude toward magic} – how locals react to spellcasters.
  \item \textbf{Typical dress / customs} – to help describe your character.
  \item \textbf{Player hook} – a story seed to tie your character to this place.
  \item \textbf{Key stakeholders} – major NPCs and factions.
  \item \textbf{Mini‑generator} – quick prompts for GMs and players (d6 or d8 table).
\end{itemize}
For GM‑only secrets (hidden agendas, faction clocks), see the \textit{GM Guide}.
\end{tcolorbox}

\section*{Quick Reference by Subgenre}
\addcontentsline{toc}{section}{Quick Reference by Subgenre}
{\small
\begin{longtable}{lll}
\toprule
\textbf{Region} & \textbf{Fantasy Subgenre} & \textbf{Magic Attitude} \\
\midrule
Acasia & Low fantasy / mercenary realism & Folk‑wards respected; visible sorcery risky. \\
Aelaerem & Folk horror / pastoral dark & Hedge‑magic common; open magic attracts the Neighbors. \\
Aeler & Industrial / survival horror & Rituals are craft; magic is guild‑regulated. \\
Aelinnel & Fae / mathematical / courtesy horror & Fae rites only; iron is taboo. \\
Black Banners & Military / mercenary warfare & Pragmatic; useful magic tolerated. \\
Ecktoria & Political intrigue / legal drama & Licensed thaumaturges; unlicensed magic is crime. \\
Linn & Viking / maritime & Weather rites are sacred; false rites are sin. \\
Mistlands & Folk horror / gothic & Bell‑magic and salt – necessary but dangerous. \\
Silkstrand & Urban / mercantile / noir & Curse‑magic feared; trade magic accepted. \\
Theona & Celtic / folk horror / taboo‑driven & Old rites respected; the Ninth taboo governs all. \\
Thepyrgos & Byzantine / academic & Magic studied; unorthodox styles unfashionable. \\
Ubral & Highland clan / honour culture & Outsider mages watched; clan shamans respected. \\
Valewood & Fairy tale / dark forest & Fey magic, imperial echoes – dangerous to meddle. \\
Vhasia & Courtly romance / fractured kingdom & Church distrusts magic; lords employ mages anyway. \\
Viterra & Legal drama / frontier justice & Hedge‑magic common; registered magic tolerated. \\
Ykrul & Steppe nomad / saga & Geometry of oaths; shamanism respected. \\
Zakov & Pirate / noir / crime syndicate & Any magic welcome – if it pays. \\
\bottomrule
\end{longtable}
}

% ============================================================
% ACASIA
% ============================================================
\section{Acasia --- ``The Broken Province''}
\index{Acasia}
\textbf{Genre:} Low fantasy / mercenary realism / border feud.

\paragraph{Attitude toward magic:} Hedge‑witches respected; visible sorcery can start a levy. The Curse is blamed on ancient magic.

\paragraph{Typical dress / customs:} Wool and leather, scarves to hide clan markings. Feasts often end in oath‑swearing; never sign a contract on a bridge.

\paragraph{Player hook:} \emph{You owe a blood‑debt to a free company captain – or you are that captain.}

\paragraph{Key stakeholders:} Margravine of the Broken March, The Lame King, Blackwood Matriarchs, Free Company Captains.

\begin{tcolorbox}[colback=gray!5!white, colframe=gray!70!black, title=Acasia Mini‑Generator (d6)]
\begin{enumerate}
\item A bridge toll has tripled overnight; the toll‑keeper is a deserter from your old company.
\item A hedge‑witch offers to break your blood‑debt – for the memory of your first kill.
\item The Lame King travels with a mute jester who writes prophecies on fallen leaves.
\item Bandits stole a Blackwood Matriarch's curse‑box. Recover it before it opens.
\item A free company captain wants to hire you – but their last three hires died of ``accidents''.
\item The Margravine's seal is missing. Whoever finds it can rewrite border law.
\end{enumerate}
\end{tcolorbox}

% ============================================================
% AELAEREM
% ============================================================
\section{Aelaerem (Halflings) --- ``Hearth \& Hollow''}
\index{Aelaerem}
\textbf{Genre:} Folk horror / pastoral quiet with dark undercurrents.

\paragraph{Attitude toward magic:} Hearth‑magic (red thread, cup‑marks) is everyday. Open magic may attract the Neighbors.

\paragraph{Typical dress / customs:} Bright wool, red doors, nine cups never poured. Hospitality is sacred – and a trap.

\paragraph{Player hook:} \emph{The Pale Shepherd came for your grandmother. You still hear her footsteps.}

\paragraph{Key stakeholders:} Apple‑Matron, The Pale Shepherd, Mummers' Captain, Lantern‑Warden.

\begin{tcolorbox}[colback=gray!5!white, colframe=gray!70!black, title=Aelaerem Mini‑Generator (d6)]
\begin{enumerate}
\item You poured the ninth cup by accident. The Hollow Bride left a silver coin under your pillow.
\item A red thread appears around your wrist each morning. Who ties it – and what do they want?
\item The Apple‑Matron's orchard produced blood‑red fruit. Eating one lets you speak with crows.
\item A mummers' troupe wears masks of your dead relatives. They invite you to join.
\item The Lantern‑Warden is deaf. She hires you to ``hear'' which bell‑line is failing.
\item The Pale Shepherd walks beside you for one night. You may ask one question – but you must answer one truth in return.
\end{enumerate}
\end{tcolorbox}

% ============================================================
% AELER
% ============================================================
\section{Aeler (Dwarves) --- ``Crowns \& Under‑Vaults''}
\index{Aeler}
\textbf{Genre:} Industrial / political / survival horror (underground).

\paragraph{Attitude toward magic:} Rituals are a craft, guild‑licensed. Unsanctioned magic damages air and stone.

\paragraph{Typical dress / customs:} Heavy aprons, breath‑masks, bell‑tokens. Count your breaths in eights; never waste water.

\paragraph{Player hook:} \emph{You broke a seal in the deep ways and something followed you up.}

\paragraph{Key stakeholders:} Vault‑Queen, High King Beneath the Peaks, True Masons, Edgewalkers.

\begin{tcolorbox}[colback=gray!5!white, colframe=gray!70!black, title=Aeler Mini‑Generator (d6)]
\begin{enumerate}
\item The ventilation system is ticking – you have one hour before the lower galleries suffocate.
\item A vein‑serpent wakes. It will trade a lost treasure for your last breath.
\item The Vault‑Queen's ledger shows a keystone that was never carved. You must invent it.
\item Edgewalkers report singing stones. The song is in a dialect that has been dead for three centuries.
\item A rival clan sabotaged your breath‑counters. Find their spy before the audit.
\item The High King Beneath the Peaks demands a new oath – carved into your own collarbone.
\end{enumerate}
\end{tcolorbox}

% ============================================================
% AELINNEL
% ============================================================
\section{Aelinnel (Gnomes) --- ``Stone, Bough, and Bright Things''}
\index{Aelinnel}
\textbf{Genre:} Fae / mathematical / courtesy horror.

\paragraph{Attitude toward magic:} Fae rites, geasa, and ledger‑magic. Iron is taboo; copper is polite.

\paragraph{Typical dress / customs:} Bright coats, copper tools, counting beads. Never trust a single ledger.

\paragraph{Player hook:} \emph{You spoke a word you shouldn't have, and now the Green Gate demands payment.}

\paragraph{Key stakeholders:} Stone Prince of Aelinnel, Lady of Thorns, Green Knight, Oath‑Carver.

\begin{tcolorbox}[colback=gray!5!white, colframe=gray!70!black, title=Aelinnel Mini‑Generator (d6)]
\begin{enumerate}
\item The Green Gate appears. It asks: ``What is the number of a promise that was never spoken?''
\item A Thorn Courtier offers you a cup of tea. If you drink, you forget your own name.
\item The Oath‑Carver has gone mad, carving the same impossible equation into every door.
\item Copper coins scream when touched. They are trying to recite a forgotten treaty.
\item The Stone Prince's throne is a living geas. Sitting on it makes you honour one ancient oath.
\item A tide‑path opens at midnight. It leads to a library where every book is blank – but whispering.
\end{enumerate}
\end{tcolorbox}

% ============================================================
% BLACK BANNERS
% ============================================================
\section{Black Banners --- ``Condotta \& Crowns''}
\index{Black Banners}
\textbf{Genre:} Military / mercenary / low fantasy warfare.

\paragraph{Attitude toward magic:} Battle‑magic is useful; unreliable mages don't get paid. Condotta captains hire cantors for morale.

\paragraph{Typical dress / customs:} Tabards over rusted mail, scarves to hide old loyalties. Payday is sacred; betrayal is business.

\paragraph{Player hook:} \emph{Your company was sold out from under you. Now you hunt the broker who took the coin.}

\paragraph{Key stakeholders:} Banner‑Captain, The Black Colonel, Clan‑Mother (Ykrul), Veteran Sergeant.

\begin{tcolorbox}[colback=gray!5!white, colframe=gray!70!black, title=Black Banners Mini‑Generator (d6)]
\begin{enumerate}
\item Payday is late. The troops are muttering. The quartermaster has disappeared.
\item The Black Colonel's banner is found torn – but no one saw it happen.
\item A rival condotta offers double pay to switch sides – but your current captain holds your hostage.
\item The Veteran Sergeant knows a secret passage into the enemy camp. He also knows it's cursed.
\item A cantor's song heals the wounded but slowly erases their memories of home.
\item The Clan‑Mother sends a war‑host. She demands a blood‑price paid in songs, not gold.
\end{enumerate}
\end{tcolorbox}

% ============================================================
% ECKTORIA
% ============================================================
\section{Ecktoria --- ``Marble \& Fire''}
\index{Ecktoria}
\textbf{Genre:} Political intrigue / legal drama / fallen empire.

\paragraph{Attitude toward magic:} Licensed thaumaturges file casting notices. Unlicensed magic is a crime, often punished by the Everflame inquisitors.

\paragraph{Typical dress / customs:} Togas, seal rings, procession masks. Law is a weapon; never refuse a toast in a coin‑house.

\paragraph{Player hook:} \emph{You forged a document that could topple a senator – or save a heretic.}

\paragraph{Key stakeholders:} Grand Magistrate, High Priest of the Everflame, Duchess‑Regent, Censor's Clerk.

\begin{tcolorbox}[colback=gray!5!white, colframe=gray!70!black, title=Ecktoria Mini‑Generator (d6)]
\begin{enumerate}
\item The Grand Magistrate's son has been cursed. The only cure is a confession that would ruin a family.
\item A Censor's clerk offers you a ``redacted'' page from the imperial archives. It shows your name.
\item The Everflame's flame gutters for one second. In that second, you see a ghost wearing your face.
\item A legal duel is declared. Your champion must fight – but the opponent's blade is made of frozen lies.
\item The Duchess‑Regent hosts a masked ball. Each mask is a binding contract.
\item A heretic claims the Everflame is not divine – it's a prison for a forgotten angel. The inquisitors are already on their way.
\end{enumerate}
\end{tcolorbox}

% ============================================================
% LINN
% ============================================================
\section{Linn --- ``Skerries \& Storm‑Oaths''}
\index{Linn}
\textbf{Genre:} Viking / maritime / oath‑bound raiding.

\paragraph{Attitude toward magic:} Weather magic is sacred; false rites are crimes against the ship. Skalds use songs for law and memory.

\paragraph{Typical dress / customs:} Wool, leather, whale‑bone runes. Oaths sworn on the tide cannot be broken without drowning.

\paragraph{Player hook:} \emph{Your father's sword was lost in a raid. The sea‑queen says the drowned still carry it.}

\paragraph{Key stakeholders:} Sea‑Queen, High Jarl, Volva of the Mist, Skald.

\begin{tcolorbox}[colback=gray!5!white, colframe=gray!70!black, title=Linn Mini‑Generator (d6)]
\begin{enumerate}
\item A drowned longship rises from the fjord. The crew is still aboard – and still owes a debt.
\item The Volva of the Mist demands you cut your shadow loose. It will serve as a spy for one year.
\item A skald sings a verse that names you as the heir to a curse you never knew existed.
\item The Sea‑Queen's crown is made of a single whale's tooth. It whispers the names of traitors.
\item A storm‑oath is broken. The wind will not return until the oath is sworn again in blood.
\item The High Jarl's funeral pyre burns for nine days. On the ninth night, he sits up and asks for his sword.
\end{enumerate}
\end{tcolorbox}

% ============================================================
% MISTLANDS
% ============================================================
\section{Mistlands --- ``Bells, Salt, and Breath''}
\index{Mistlands}
\textbf{Genre:} Folk horror / gothic / supernatural dread.

\paragraph{Attitude toward magic:} Bell‑lines and salt are the only trusted magic. Outsider spellcasters are viewed with suspicion.

\paragraph{Typical dress / customs:} Grey wool, fog‑veils, bell‑charms. Never breathe deeply after the ninth bell.

\paragraph{Player hook:} \emph{You were born during a bell‑line failure. The fog knows your name – it has visited twice.}

\paragraph{Key stakeholders:} Legate of the Mists, Bell‑Warden, Salt‑Monk, Mist‑Seer.

\begin{tcolorbox}[colback=gray!5!white, colframe=gray!70!black, title=Mistlands Mini‑Generator (d6)]
\begin{enumerate}
\item A bell‑line goes silent. The fog pours through the gap, carrying a voice that asks for a story.
\item The Salt‑Monk's salt pans have turned black. The curse will spread unless you find a new vein.
\item A Mist‑Seer offers to read your future in the fog. The price is one memory you never share.
\item The Legate of the Mists is an Aeler exile. He will trade safe passage for a secret about your past.
\item A corpse on the bell‑line has no face – but it wears your grandmother's wedding ring.
\item The fog parts for one minute, revealing a city that should not exist. Its bell is ringing the ninth time.
\end{enumerate}
\end{tcolorbox}

% ============================================================
% SILKSTRAND
% ============================================================
\section{Silkstrand --- ``City of Bridges \& Dyewater''}
\index{Silkstrand}
\textbf{Genre:} Urban / mercantile / noir.

\paragraph{Attitude toward magic:} Curse‑magic feared; trade magic (dye‑wards, bridge‑blessings) accepted. The Matron employs a weft‑reader.

\paragraph{Typical dress / customs:} Dyed silks, debtor's stains, bridge tokens. Every bridge is a ledger; every handshake a contract.

\paragraph{Player hook:} \emph{You're the only one who knows the secret of the red dye. The guild wants the recipe – or your hands.}

\paragraph{Key stakeholders:} The Matron, Guildmistress of Dyers, Bridge Bailiff, Night‑Boat Smuggler (``Ravel'').

\begin{tcolorbox}[colback=gray!5!white, colframe=gray!70!black, title=Silkstrand Mini‑Generator (d6)]
\begin{enumerate}
\item A bolt of cloth bleeds when cut. The blood shows a murder that hasn't happened yet.
\item The Bridge Bailiff auctions a debt. You can buy it – but the debtor is a ghost.
\item The Night‑Boat Smuggler offers you a package. Do not open it. Do not ask what it is.
\item The Guildmistress of Dyers demands a new colour – dyed from the ashes of a broken oath.
\item The Matron's weft‑reader sees a thread connecting you to the Drowned Man. She refuses to cut it.
\item A dye‑vat overflows, turning a canal to ink. Words form on the surface – a list of names. Yours is last.
\end{enumerate}
\end{tcolorbox}

% ============================================================
% THEONA
% ============================================================
\section{Theona --- ``Three Greens, No Ninth''}
\index{Theona}
\textbf{Genre:} Celtic / folk horror / taboo‑driven.

\paragraph{Attitude toward magic:} Old rites (well‑blessings, bride‑peace) are essential. Breaking a taboo invites the Green Host.

\paragraph{Typical dress / customs:} Peat‑stained cloaks, hawthorn pins, counting in eights. The ninth cup is never poured.

\paragraph{Player hook:} \emph{You were at a bride‑peace feast when the ninth name was spoken. The Green Host rode – but you survived.}

\paragraph{Key stakeholders:} Matron of Wells, Three‑Isles King, Lady Beneath, Taboo‑Witness.

\begin{tcolorbox}[colback=gray!5!white, colframe=gray!70!black, title=Theona Mini‑Generator (d6)]
\begin{enumerate}
\item A well runs dry. The Matron of Wells says the water was stolen – by your own reflection.
\item The Three‑Isles King marries a second time. The bride is a woman you killed (or failed to save).
\item The Taboo‑Witness catches you pouring the ninth cup. She offers a deal: your right hand, or a memory.
\item The Green Host rides tonight. You can join them – but you will forget how to return.
\item The Lady Beneath sends a dream: a coin lies at the bottom of the marsh. Take it, and she will remember your name.
\item A bride‑peace is broken. The groom's family demands a blood‑price; the bride's family offers you instead.
\end{enumerate}
\end{tcolorbox}

% ============================================================
% THEPYRGOS
% ============================================================
\section{Thepyrgos --- ``City of a Thousand Stairs''}
\index{Thepyrgos}
\textbf{Genre:} Byzantine / academic / arcane mystery.

\paragraph{Attitude toward magic:} The Synod licenses magic; unorthodox casting is not illegal – merely unfashionable. Witch‑hunters conscript talent.

\paragraph{Typical dress / customs:} Academic robes, bell‑badges, stair‑passes. Every stair climbed is a word in a conversation.

\paragraph{Player hook:} \emph{You tested as ``gifted'' on the candle‑smoke test. The university wants you; the Chain‑Lanterns want you in chains.}

\paragraph{Key stakeholders:} The Archon, Lighthouse‑Patriarch, Nomophylax, Palkar Captain.

\begin{tcolorbox}[colback=gray!5!white, colframe=gray!70!black, title=Thepyrgos Mini‑Generator (d6)]
\begin{enumerate}
\item The Lighthouse‑Patriarch has not spoken in three years. His shadow gives lectures on forbidden texts.
\item A stair token is stolen. Without it, you cannot reach the upper archives – and a demon is loose there.
\item The Nomophylax declares your magical talent ``unlicensed''. You have one day to pass the Synod exam.
\item A Palkar Captain offers to erase your criminal record – in exchange for a ritual performed in the cisterns.
\item The Archon dies. His last word is your name. The election for his successor begins at midnight.
\item A student has been erased from every record. Only you remember them. The Chain‑Lanterns want to know why.
\end{enumerate}
\end{tcolorbox}

% ============================================================
% UBRAL
% ============================================================
\section{Ubral --- ``The Stone Between Spears''}
\index{Ubral}
\textbf{Genre:} Highland clan / reiver / honour culture.

\paragraph{Attitude toward magic:} Clan shamans are respected; outsider mages are watched. Stone‑speakers (dwarf‑trained) are tolerated.

\paragraph{Typical dress / customs:} Tweed, iron brooches, guest‑tokens. Never refuse a cup from a hearth‑aunt – hospitality is a contract.

\paragraph{Player hook:} \emph{You carry a wergild debt that your father couldn't pay. The feud is now yours.}

\paragraph{Key stakeholders:} Lady of Tor, Council of Cairns, Feud‑Broker, Wergild Counter.

\begin{tcolorbox}[colback=gray!5!white, colframe=gray!70!black, title=Ubral Mini‑Generator (d6)]
\begin{enumerate}
\item The Council of Cairns adds a new stone. It has your name carved into it – but you are not dead.
\item A guest‑token is broken. You are now an outlaw unless you find the piece that was stolen.
\item The Feud‑Broker offers to settle your wergild – but the payment is the hand of a loved one.
\item The Lady of Tor's son was killed by a monster. The monster is actually a cursed ancestor of yours.
\item A Wergild Counter admits he has miscounted. The debt is actually double. He wants you to kill the witness.
\item A stone‑speaker claims the mountain is waking. It demands a tribute of fresh blood.
\end{enumerate}
\end{tcolorbox}

% ============================================================
% VALEWOOD
% ============================================================
\section{Valewood --- ``Empire Under Leaves''}
\index{Valewood}
\textbf{Genre:} Fairy tale / dark forest / fey remnant.

\paragraph{Attitude toward magic:} Imperial echoes and fae rites are dangerous. Never eat fruit from a tree with a face.

\paragraph{Typical dress / customs:} Leather, moss‑dyed cloaks, way‑cords. Speak in riddles or not at all.

\paragraph{Player hook:} \emph{A star‑road opened in your dream. You walked it – now something walks behind you.}

\paragraph{Key stakeholders:} Hazel Queen, Alder King, The Huntsman, Pathweaver.

\begin{tcolorbox}[colback=gray!5!white, colframe=gray!70!black, title=Valewood Mini‑Generator (d6)]
\begin{enumerate}
\item A star‑road appears at midnight. It leads to a throne made of antlers. The Alder King sits on it – and he expects you to kneel.
\item The Pathweaver offers you a cord. If you tie it, you will never be lost again. If you untie it, you will forget where home is.
\item The Hazel Queen is dying. Her last request: find the acorn of the first tree, buried under a court that no longer exists.
\item The Huntsman's hounds are loose. They will not attack – they will only follow you until you speak your true name.
\item An imperial echo (a ghost of a long‑dead soldier) mistakes you for an enemy. It will not listen to reason – only to a duellist's blade.
\item A tree with a face offers you its fruit. Eat it, and you will understand the language of stone. Refuse, and it will uproot itself and follow you.
\end{enumerate}
\end{tcolorbox}

% ============================================================
% VHASIA
% ============================================================
\section{Vhasia --- ``The Fractured Sun''}
\index{Vhasia}
\textbf{Genre:} Courtly romance / chivalric / fractured kingdom.

\paragraph{Attitude toward magic:} The Church of the Flame distrusts arcane magic; secular lords employ court mages anyway. Chevauchée raiders use cantors for morale.

\paragraph{Typical dress / customs:} Velvet, partisan ribbons, duelling scars. Never carry a coin with only one sunburst – treason.

\paragraph{Player hook:} \emph{You're the Last Dauphin – or a convincing fake. Either way, armies march for your claim.}

\paragraph{Key stakeholders:} Two Crowns, Queen‑Mother, Last Dauphin, Marshal‑Exile.

\begin{tcolorbox}[colback=gray!5!white, colframe=gray!70!black, title=Vhasia Mini‑Generator (d6)]
\begin{enumerate}
\item The Queen‑Mother offers you a marriage alliance. The bride is a ghost bound to a tapestry.
\item The Marshal‑Exile demands a duel. The winner takes the army; the loser takes a beheading.
\item A chevauchée burns a village that swore neutrality. The only survivor is a child who calls you by a dead knight's name.
\item The Two Crowns are secretly the same person – a master of disguise who enjoys war.
\item The Last Dauphin is a commoner recruited by a desperate lord. He does not know he is a fake.
\item A duelling scar heals overnight, leaving silver thread in your skin. The wound now whispers treason.
\end{enumerate}
\end{tcolorbox}

% ============================================================
% VITERRA
% ============================================================
\section{Viterra --- ``The Hedge‑Law Realm''}
\index{Viterra}
\textbf{Genre:} Legal drama / low fantasy / frontier justice.

\paragraph{Attitude toward magic:} Hedge‑magic is common and unregulated; any thaumaturgy that leaves a trace must be registered. Legal duels may involve wizards.

\paragraph{Typical dress / customs:} Wool, seal rings, boundary stones. Every hedge is a precedent; every dike a court ruling.

\paragraph{Player hook:} \emph{Your family lost the border dispute. You have one year to find a lost charter or lose the land forever.}

\paragraph{Key stakeholders:} Warrior Queen, Queen's Justiciar, Fen Reeve, Master Advocate Halric.

\begin{tcolorbox}[colback=gray!5!white, colframe=gray!70!black, title=Viterra Mini‑Generator (d6)]
\begin{enumerate}
\item The Warrior Queen comes in person to settle a dispute. She asks each party one question. The liar loses a hand.
\item The Fen Reeve's tally‑stick is broken. Until it is mended, the river toll is doubled – and the river itself is angry.
\item Master Advocate Halric offers to defend you for free. His last three clients all died mysteriously.
\item A boundary stone moves three feet east every night. The neighbours swear they are not touching it.
\item The Queen's Justiciar arrests you for ``hedge‑walking without a license''. The sentence is a legal duel at dawn.
\item A charter written in disappearing ink is the only proof of your inheritance. The ink will be gone by sunrise.
\end{enumerate}
\end{tcolorbox}

% ============================================================
% YKRUL
% ============================================================
\section{Ykrul (Orcs) --- ``Wolf Standards, Winter Camps''}
\index{Ykrul}
\textbf{Genre:} Steppe nomad / saga / honour culture.

\paragraph{Attitude toward magic:} Geometry of oaths (Kon'reh) and sky‑speakers. Written magic is suspicious; sung magic is respected.

\paragraph{Typical dress / customs:} Leather, horse‑hair braids, wolf‑tooth standards. Never ride between two wolves at dusk.

\paragraph{Player hook:} \emph{You broke a hostage string. Now the Khatun demands a blood‑price – or your youngest kin.}

\paragraph{Key stakeholders:} Khatun of Ring, Khagan's Kin, Sky‑Speaker, Kon'reh Masters.

\begin{tcolorbox}[colback=gray!5!white, colframe=gray!70!black, title=Ykrul Mini‑Generator (d6)]
\begin{enumerate}
\item A white squall separates you from the herd. In the snow, you see the faces of everyone you have wronged.
\item The Sky‑Speaker calls a storm to halt an enemy army. The storm demands a sacrifice: your left ear.
\item A Kon'reh Master sets you a problem: ``How many camps can be pitched before the river freezes?'' The wrong answer is death.
\item The Khagan's Kin offers you a hostage string. If you accept, you gain a loyal servant – but also a rival who will hunt you.
\item A wolf‑standard is stolen. The clan will not fight until it is returned. The thief is a ghost who died in the last war.
\item The Khatun of Ring declares a blood‑price on your head. You have one day to pay or one day to ride.
\end{enumerate}
\end{tcolorbox}

% ============================================================
% ZAKOV
% ============================================================
\section{Zakov --- ``Salt \& Serpent''}
\index{Zakov}
\textbf{Genre:} Pirate / noir / crime syndicate.

\paragraph{Attitude toward magic:} Any magic is welcome – if it pays. Cursed cargoes are a common hazard; fae bargains are made at the Crimson Docks.

\paragraph{Typical dress / customs:} Brine‑stained coats, reef‑tooth necklaces, smuggler's tokens. Never take a coin that tastes of brine – the Serpent's Mark is a debt.

\paragraph{Player hook:} \emph{You have a syndicate marker that could buy a fleet – or get you killed. The Salt Prince wants to collect.}

\paragraph{Key stakeholders:} Salt Prince, Pirate Queen, Kraken's Tongue, Syndicate Silent.

\begin{tcolorbox}[colback=gray!5!white, colframe=gray!70!black, title=Zakov Mini‑Generator (d6)]
\begin{enumerate}
\item The Kraken's Tongue offers a prophecy. To hear it, you must drink a cup of seawater that tastes like your mother's tears.
\item A Pirate Queen's crew is cursed: they cannot step on land until they pay a debt to a man who drowned fifty years ago.
\item The Salt Prince's collector comes for your marker. You can pay with a memory, a finger, or a voyage to a reef that does not exist.
\item A fae bargain is struck on the Crimson Docks. The price: a stolen hour of sleep from every sailor in port.
\item A smuggler's token is found in your pocket. It belonged to someone who was hanged last week. The guard assumes you are the accomplice.
\item A reef‑tooth necklace bleeds when you touch it. It belonged to a pirate queen who was betrayed – and she is not dead.
\end{enumerate}
\end{tcolorbox}

% ============================================================
% OTHER REGIONS (BRIEF)
% ============================================================
\section{Other Regions (Brief)}
\label{sec:cultures-other}
For completeness, these areas appear in the full setting but are less central to the core game. Ask your GM if your character may hail from one of them.
\begin{itemize}
  \item \textbf{Dungeon – ``Living Infrastructure''} – Mythic underworld, survival horror.
  \item \textbf{The Ways Between – ``Spiritways \& Veilways''} – Liminal paths, dream‑logic.
  \item \textbf{The Wilds – ``Roads, Ruins, and Weather''} – A reskin palette for any terrain.
  \item \textbf{Vilikari – ``Laurels \& Longhouses''} – Federated frontier, two laws.
\end{itemize}

% ============================================================
% PEOPLES OF FATE'S EDGE — DETAILED MECHANICS
% ============================================================
\section{Peoples of Fate's Edge — Detailed Mechanics}
\label{sec:peoples-mechanics}

The following sections provide the mechanical traits (attribute adjustments, skill bonuses, cultural talents, Strings, sample Bonds, and sample Complications) for the major ancestries of the Amaranthine. Use these during character creation to bring your character's heritage to life.

% ------------------------------------------------------------
% SMALL FOLK OF THE THRESHOLD (Aelaerem & Aelinnel)
% ------------------------------------------------------------
\subsection{Small Folk of the Threshold}
\label{sec:small-folk-detailed}
\index{Aelaerem}\index{Aelinnel}\index{Small Folk}

The Aelaerem (halflings) and Aelinnel (gnomes) are diminutive peoples attuned to liminal spaces – doorways, hedge‑gaps, tide‑cut stairs, and the quiet places between hearth and hollow.

\paragraph{Restriction:} Cannot use Heavy Armor. May use any one‑handed or two‑handed weapon, but Heavy Weapons impose \textbf{–1 Position} on all actions while wielding them (before any other Position modifiers).

\paragraph{Bonus:} When wielding a \textbf{spear, quarterstaff, or light polearm} (any Medium or Light weapon with the \texttt{Reach} tag), treat the weapon as \textbf{one weight class lighter} for the purpose of attack modifiers. Additionally, once per scene when you make a melee attack with such a weapon, you may spend \textbf{1 Boon} to convert a \textit{Partial} success into a \textit{Clean Success} (representing a clever thrust that uses the enemy's momentum against them).

\subsubsection{Aelaerem — The Hearth's Luck}
\label{sec:aelaerem-racial-detailed}
\textbf{Bonus — Once per scene}, you may improve your \textbf{Position by one step} (e.g., Desperate → Controlled, Controlled → Dominant) as a free action \textbf{before} rolling an action. This represents a flash of luck, a well‑timed sidestep, or an unexpected opening.

\begin{tcolorbox}[colback=green!5!white,colframe=green!70!black,title=From the Hearth Ledger]
\emph{``The Aelaerem say that a door left unlatched is an invitation to the Hollow. But sometimes, an unlatched door is just an escape.''}  
Use your luck wisely. The Hollow is patient.
\end{tcolorbox}

\subsubsection{Aelinnel — The Bright Trick}
\label{sec:aelinnel-racial-detailed}
Choose \textbf{one} of the following abilities:

\begin{description}
  \item[Option A: Short Step (Teleport)] \hfill \\
  \textbf{Bonus — Once per scene}, as a \textbf{Move action}, you may instantly teleport to any location you can see within \textbf{Near} range that is not currently observed by any enemy. You cannot pass through solid barriers. This is a fey‑like blink, not a physical jump.
  
  \item[Option B: Knack for the Right Thing (Handy Item)] \hfill \\
  \textbf{Bonus — Once per scene}, you may declare that you have a \textbf{small, mundane item} on your person that is useful in the current situation (e.g., a coil of rope, a lockpick, a candle, a chalk stick). This costs no Boon and requires no roll. The item must be plausible for your character to carry.
\end{description}

\begin{tcolorbox}[colback=cyan!5!white,colframe=cyan!70!black,title=A Gnomish Maxim]
\emph{``Count the steps in odd numbers. Even is for the dead. … Wait, did I say odd? Let me check my abacus.''}  
Aelinnel never trust a single ledger. Neither should you.
\end{tcolorbox}

% ------------------------------------------------------------
% AELER (DWARVES)
% ------------------------------------------------------------
\subsection{Aeler — People of Stone, Breath, and Ledger (Dwarves)}
\label{sec:aeler-detailed}
\index{Aeler!mechanics}

\emph{``The mountain does not forget a misplaced keystone. It waits – then settles the debt.''}

\begin{tcolorbox}[colback=black!3!white,colframe=black!40!white,title=Aeler at a Glance]
\textbf{Attribute Adjustments:} Body +1, Spirit +1, Presence –1 \\
\textbf{Skill Bonuses:} Craft +1, Endurance +1, Mechanics +1 \\
\textbf{Signature Talents:} \emph{Spirit Shield} • \emph{Keystone Ledger} \\
\textbf{Strings:} Charter‑Stamped, Stone‑Speech Initiate
\end{tcolorbox}

\paragraph{Cultural Talents}

\textbf{Minor (2–3 XP)}
\begin{itemize}
  \item \textbf{Stone‑Sense (2 XP):} Passive. +1 die to detect flaws in stone, hidden passages, or structural weaknesses underground.
  \item \textbf{Breath‑Counter (2 XP):} Active, once per scene. When in a sealed space, reduce Fatigue from lack of air by 1 level.
  \item \textbf{Oath‑Cord (3 XP):} Active, once per session. Tie a cord to bind a simple promise; breaking it costs the violator 1 SB (Hearts).
\end{itemize}

\textbf{Major (4–6 XP)}
\begin{itemize}
  \item \textbf{Spirit Shield (6 XP):} Passive. When you Defend, you may convert the first Harm 1 you would take into Fatigue (minimum 1 Fatigue). \textit{Requires: Spirit 3+.}
  \item \textbf{Keystone Ledger (4 XP):} Passive. You may keep a written record of one Debt Clock without expending an action to recall it.
\end{itemize}

\textbf{Prestige (7–10 XP)}
\begin{itemize}
  \item \textbf{Forge‑Patriarch (10 XP):} Active, once per campaign. Establish a \textbf{forge‑citadel} (Major Asset) and gain a loyal cadre of smiths and engineers (Cap 3 followers). This requires a season of downtime and significant resources.
\end{itemize}

\paragraph{Sample Bonds}
\begin{itemize}
  \item ``I owe my life to the Vault‑Queen's wardens.''
  \item ``My clan's keystone is lost – I must find it before the mountain settles.''
\end{itemize}

\paragraph{Sample Complications}
\begin{itemize}
  \item A sealed hold refuses you entry.
  \item A rival clan claims your keystone.
  \item A White Flood drowns your route.
\end{itemize}

% ------------------------------------------------------------
% AELINNEL (GNOMES)
% ------------------------------------------------------------
\subsection{Aelinnel — People of Sums, Bough, and Bright Things (Gnomes)}
\label{sec:aelinnel-detailed}
\index{Aelinnel!mechanics}

\emph{``Count the steps in odd numbers. Even is for the dead.''}

\begin{tcolorbox}[colback=black!3!white,colframe=black!40!white,title=Aelinnel at a Glance]
\textbf{Attribute Adjustments:} Wits +1, Spirit +1, Body –1 \\
\textbf{Skill Bonuses:} Investigation +1, Lore +1, Mechanics +1 \\
\textbf{Signature Talents:} \emph{Oath‑Carver} • \emph{Green Gate Key} \\
\textbf{Strings:} Hall‑Stamped, Tide‑Path Known
\end{tcolorbox}

\paragraph{Cultural Talents}
\textit{Note: The \emph{Short Step} and \emph{Knack for the Right Thing} are part of the Small Folk racial package (see above), not purchased talents.}

\textbf{Minor (2–3 XP)}
\begin{itemize}
  \item \textbf{Counting Etiquette (2 XP):} Passive. Once per scene, when you carefully count (steps, breaths, beads), you may shift Position +1 for a single action requiring rhythm or precision.
  \item \textbf{Copper over Iron (2 XP):} Passive. When you visibly favor copper or brass tools over iron in fae‑facing scenes, you gain +1 die on Courtesy or Bargain rolls.
\end{itemize}

\textbf{Major (4–6 XP)}
\begin{itemize}
  \item \textbf{Oath‑Carver (4 XP):} Active. Set a promise in quartz (a visible physical token). The target gains +1 die to keep the promise; breaking it costs 1 SB (Spades). You may maintain one quartz oath at a time.
\end{itemize}

\textbf{Prestige (7–10 XP)}
\begin{itemize}
  \item \textbf{Green Gate Key (8 XP):} Passive. Once per journey, you may declare a hidden tide‑path or stone‑door that halves a travel leg (reduce Travel Clock by half its segments). The GM gains 1 SB as the path demands a small toll (a memory, a name).
\end{itemize}

\paragraph{Sample Bonds}
\begin{itemize}
  \item ``A Thorn Courtier saved me from a geas – now I owe them a favor.''
  \item ``My grandmother taught me the ninth silence.''
\end{itemize}

\paragraph{Sample Complications}
\begin{itemize}
  \item A geas catches you in a careless word.
  \item The Green Gate opens at the wrong hour.
  \item A Thorn Court summons you.
\end{itemize}

% ------------------------------------------------------------
% AELAEREM (HALFLINGS)
% ------------------------------------------------------------
\subsection{Aelaerem — People of Hearth \& Hollow (Halflings)}
\label{sec:aelaerem-detailed}
\index{Aelaerem!mechanics}

\emph{``A door left unlatched is an invitation to the Hollow. But sometimes, an unlatched door is just an escape.''}

\begin{tcolorbox}[colback=black!3!white,colframe=black!40!white,title=Aelaerem at a Glance]
\textbf{Attribute Adjustments:} Wits +1, Presence +1, Body –1 \\
\textbf{Skill Bonuses:} Survival +1, Sway +1, Stealth +1 \\
\textbf{Signature Talents:} \emph{Barrow‑Wise} • \emph{Pale Shepherd's Grace} \\
\textbf{Strings:} Red‑Thread Guest, Barrow‑Friend
\end{tcolorbox}

\paragraph{Cultural Talents}
\textit{Note: The \emph{Hearth's Luck} ability is part of the Small Folk racial package (see above), not a purchased talent.}

\textbf{Minor (2–3 XP)}
\begin{itemize}
  \item \textbf{Red Thread (2 XP):} Active, once per session. Tie a red thread on a door or bridge; the next person who crosses must pass a Resolve test (DV 2) or forget why they entered.
  \item \textbf{Cup‑Mark (3 XP):} Active. Leave a cup‑mark on a stone; when you return before dawn, you may ignore one minor social complication (e.g., a suspicious guard).
\end{itemize}

\textbf{Major (4–6 XP)}
\begin{itemize}
  \item \textbf{Barrow‑Wise (4 XP):} Passive. +1 die to Lore about the Neighbors, the Hollow, or ancient barrows. Once per session, you may ask the GM one yes/no question about a supernatural threat in the current scene.
  \item \textbf{Mummers' Mask (6 XP):} Active, once per scene. Don a mask and adopt a role; for the next exchange, you may use \emph{Presence + Performance} in place of any social skill.
\end{itemize}

\textbf{Prestige (7–10 XP)}
\begin{itemize}
  \item \textbf{Pale Shepherd's Grace (8 XP):} Passive. When you would take Harm, you may instead mark that amount as Fatigue. If you do, the GM gains 1 SB (Hearts) – the Hollow watches you. \textit{Requires: Spirit 3+.}
\end{itemize}

\paragraph{Sample Bonds}
\begin{itemize}
  \item ``The Apple‑Matron saved my family from famine – I owe her a harvest.''
  \item ``The Pale Shepherd walked past me once. I still don't know why.''
\end{itemize}

\paragraph{Sample Complications}
\begin{itemize}
  \item You pour the ninth cup; the Hollow answers.
  \item A mummers' mask sticks to your face.
\end{itemize}

% ------------------------------------------------------------
% LETHAI (ELVES)
% ------------------------------------------------------------
\subsection{The Lethai — Root, River, \& Roof‑Tree / Mind's Eye \& Civic Measure}
\label{sec:lethai-detailed}
\index{Lethai}

\emph{``Two paths, one blood – the Sundered Kin.''}

\subsubsection{Lethai-al (Wood Elves)}
\label{subsec:lethai-al-detailed}
\index{Wood Elves!mechanics}

\begin{tcolorbox}[colback=black!3!white,colframe=black!40!white,title=Wood Elf at a Glance]
\textbf{Attribute Adjustments:} Body +1, Wits +1, Presence –1 \\
\textbf{Skill Bonuses:} Survival +1, Stealth +1, Nature +1 \\
\textbf{Signature Talents:} \emph{Root‑Law} • \emph{Green Ward} \\
\textbf{Strings:} Root‑Oath Bound, Tree‑Friend
\end{tcolorbox}

\paragraph{Cultural Talents}

\textbf{Minor (2–3 XP)}
\begin{itemize}
  \item \textbf{Root‑Law (2 XP):} Passive. +1 die to negotiate treaties or oaths that involve land, forest, or seasonal cycles.
  \item \textbf{Tree‑Speak (3 XP):} Active, once per scene. Commune with a single tree or plant; gain +2 dice on Investigation or Survival to understand events that happened near it.
\end{itemize}

\textbf{Major (4–6 XP)}
\begin{itemize}
  \item \textbf{Green Ward (4 XP):} Passive. When in forest terrain, you are always treated as having partial cover against ranged attacks.
  \item \textbf{Coppice Memory (6 XP):} Active, once per session. Perfectly recall a scent, sound, or footprint from up to a week ago. \textit{Requires: Wits 3+.}
\end{itemize}

\textbf{Prestige (7–10 XP)}
\begin{itemize}
  \item \textbf{Wild Speaker (10 XP):} Active, once per scene. Command a single animal or plant creature (Cap 2) to perform a simple action (flee, attack, block). The GM may award 1 SB if the command contradicts the creature's nature.
\end{itemize}

\subsubsection{Lethai-thora (High Elves)}
\label{subsec:lethai-thora-detailed}
\index{High Elves!mechanics}

\begin{tcolorbox}[colback=black!3!white,colframe=black!40!white,title=High Elf at a Glance]
\textbf{Attribute Adjustments:} Wits +1, Spirit +1, Body –1 \\
\textbf{Skill Bonuses:} Lore +1, Arcana +1, Investigation +1 \\
\textbf{Signature Talents:} \emph{Lorekeeper} • \emph{Ways‑Walker's Step} \\
\textbf{Strings:} Archive‑Key, Debate‑Sealed
\end{tcolorbox}

\paragraph{Cultural Talents}

\textbf{Minor (2–3 XP)}
\begin{itemize}
  \item \textbf{Lorekeeper (2 XP):} Passive. +1 die to recall historical or magical information. Once per session, you may ask the GM a yes/no question about a forgotten fact.
  \item \textbf{Weave Anchor (3 XP):} Passive. Reduce magical Backlash from a single 1 by one step (Minor → none, Major → Minor) once per scene.
\end{itemize}

\textbf{Major (4–6 XP)}
\begin{itemize}
  \item \textbf{Academic Immunity (4 XP):} Passive. When in a Thepyrgos library or university, you are treated as a \emph{guest scholar} – guards will not arrest you without a writ from the Archon.
  \item \textbf{Mind Over Matter (6 XP):} Active, once per scene. Re‑roll a failed Resolve or Lore test; you must accept the second roll.
\end{itemize}

\textbf{Prestige (7–10 XP)}
\begin{itemize}
  \item \textbf{Ways‑Walker's Step (10 XP):} Passive. Once per arc, you may turn a complication from a ritual or magical research into a boon (e.g., a backlash reveals a secret path). The GM banks 2 SB.
\end{itemize}

\subsubsection{Lethai-ar (Dark Elves — The Oathbound)}
\label{subsec:lethai-ar-detailed}
\index{Lethai-ar}

\emph{``Vows are blades. Wear them carefully.''}

\begin{tcolorbox}[colback=black!3!white,colframe=black!40!white,title=Dark Elf at a Glance]
\textbf{Attribute Adjustments:} Wits +1, Presence +1, Spirit –1 \\
\textbf{Skill Bonuses:} Subterfuge +1, Deception +1, Command +1 \\
\textbf{Signature Talents:} \emph{Serpent's Shed} • \emph{Spider's Mercy} \\
\textbf{Strings:} Mask‑Circle Known, Bound to a Threshold
\end{tcolorbox}

\paragraph{Cultural Talents}

\textbf{Minor (2–3 XP)}
\begin{itemize}
  \item \textbf{Mask‑Right (2 XP):} Passive. When you wear a mask (even a simple half‑mask), you gain +1 die to resist social pressure or intimidation.
  \item \textbf{Vow‑Touch (3 XP):} Active, once per scene. Touch a person and whisper a short geas (``Speak only truth''); if they fail a Resolve test (DV 3), they suffer –1 die to resist that command for one exchange.
\end{itemize}

\textbf{Major (4–6 XP)}
\begin{itemize}
  \item \textbf{Serpent's Shed (4 XP):} Active, once per session. Mark 1 Fatigue to negate one ongoing Condition (Fear, Poison, Charm) – you ``shed'' it like a skin.
  \item \textbf{Spider's Mercy (6 XP):} Active, once per scene. When you Defend and successfully protect an ally, you may transfer 1 Harm from that ally to yourself (converted to Fatigue if you have Light armor or less).
\end{itemize}

\textbf{Prestige (7–10 XP)}
\begin{itemize}
  \item \textbf{Oath‑Eater (8 XP):} Active, once per campaign. Consume a broken oath (yours or another's) to regain 2 Boons and clear 2 Fatigue. The GM gains a permanent \emph{Oath Debt} (1 SB per session) – a rival will collect.
\end{itemize}

\paragraph{Sample Bonds (all Lethai)}
\begin{itemize}
  \item ``A Lethai‑ar saved me from a wyrm – they demand a vow of silence in return.''
  \item ``My family's grove was poisoned; I carry an acorn from the mother tree.''
\end{itemize}

\paragraph{Sample Complications (all Lethai)}
\begin{itemize}
  \item You break a root‑law.
  \item A Lethai‑ar mask‑circle demands your presence.
  \item A High Elf rival publishes a refutation of your thesis.
\end{itemize}

% ------------------------------------------------------------
% YKRUL (ORCS)
% ------------------------------------------------------------
\subsection{Ykrul — The People of the Violet Steppe (Orcs)}
\label{sec:ykrul-detailed}
\index{Ykrul!mechanics}

\begin{tcolorbox}[colback=red!5!white,colframe=red!75!black,title=Note on Naming]
\textbf{``Orc'' is a slur.} The Ykrul are a proud people of the steppe; the word ``orc'' was coined by Utaran legions to dehumanize them. Wise travellers use \emph{Ykrul} or \emph{Steppe Folk}. Using ``orc'' in conversation with a Ykrul may be taken as a grave insult.
\end{tcolorbox}

\emph{``The wind carries names. Ride straight, or you will be forgotten.''}

\begin{tcolorbox}[colback=black!3!white,colframe=black!40!white,title=Ykrul at a Glance]
\textbf{Attribute Adjustments:} Body +1, Spirit +1, Presence –1 \\
\textbf{Skill Bonuses:} Riding +1, Survival +1, Command +1 \\
\textbf{Signature Talents:} \emph{Blood Memory} • \emph{Sky‑Speaker's Whistle} \\
\textbf{Strings:} Wind‑Knot Token, Banner Hospitality
\end{tcolorbox}

\paragraph{Cultural Talents}

\textbf{Minor (2–3 XP)}
\begin{itemize}
  \item \textbf{Blood Frenzy (2 XP):} Reactive. When you take Harm, you may gain +1 die on your next melee attack.
  \item \textbf{Hostage String (2 XP):} Active, once per session. Tie a cord to a willing ally; until the cord is cut, you may transfer 1 Fatigue from them to yourself as a free action.
\end{itemize}

\textbf{Major (4–6 XP)}
\begin{itemize}
  \item \textbf{Blood Memory (4 XP):} Passive. After a battle, you may gain +1 die on a single roll that mimics a foe's tactic (e.g., using a captured weapon). This lasts for the next scene.
  \item \textbf{Sky‑Speaker's Whistle (6 XP):} Active, once per scene. Imitate a steppe omen (eagle cry, thunder boom). Enemies within Near must pass a Resolve test (DV 2) or lose their next minor action.
\end{itemize}

\textbf{Prestige (7–10 XP)}
\begin{itemize}
  \item \textbf{Warglord (10 XP):} Active, once per campaign. Unite scattered warbands under a single banner. Gain a \textbf{War Host} (Major Asset) and a loyal cavalry unit (Cap 5 followers). Start \emph{Logistics} and \emph{Grudge} clocks.
\end{itemize}

\paragraph{Sample Bonds}
\begin{itemize}
  \item ``The Khagan's kin saved my life – I owe a foal each spring.''
  \item ``My hostage string was cut; I must find the betrayer before the clan does.''
\end{itemize}

\paragraph{Sample Complications}
\begin{itemize}
  \item A white squall separates you from the herd.
  \item A rival clan demands a blood‑price.
  \item The Sky‑Speaker claims you owe a storm.
\end{itemize}

% ------------------------------------------------------------
% MIXED HERITAGE
% ------------------------------------------------------------
\subsection{Mixed Heritage — Half-Elves, Half-Ykrul, Half-Others}
\label{sec:mixed-heritage}
\index{Mixed Heritage}

\emph{``You walk two roads. Both watch you.''}

\begin{tcolorbox}[colback=black!3!white,colframe=black!40!white,title=Mixed Heritage at a Glance]
\textbf{Choose one +1 and one –1 from parent cultures} \\
\textbf{Choose two skill bonuses} \\
\textbf{Access both talent lists} (subject to prerequisites) \\
\textbf{Choose primary String}
\end{tcolorbox}

\paragraph{Attribute Adjustments:} Choose \textbf{one} +1 from either parent culture, and \textbf{one} –1 from the other (or a standard +1, +1, –1 as determined with GM).

\paragraph{Skill Bonuses:} Choose \textbf{two} skill bonuses from either parent culture (cannot take the same skill twice).

\paragraph{Cultural Talents:} You may select from \textbf{both} parent cultures' talent lists, but you must meet any prerequisites (e.g., Spirit 3+ for some talents). You may not take mutually exclusive talents.

\paragraph{String Access:} You have access to the \textbf{Strings} of both cultures, but you must choose one as your primary for purposes of social perception (e.g., a half‑Aeler, half‑Lethai might be \emph{Charter‑Stamped} but not \emph{Root‑Oath Bound} unless you prove yourself).

\paragraph{Sample Bonds}
\begin{itemize}
  \item ``My father's clan disowned me; my mother's people took me in – I carry both names.''
  \item ``I am the living treaty between two feuding houses.''
\end{itemize}

\paragraph{Sample Complications}
\begin{itemize}
  \item Neither culture fully trusts you.
  \item You are called to honour two conflicting oaths.
  \item A purist from one side challenges your right to claim their traditions.
\end{itemize}